% !PRINTED 187
\weberpage{187}

Now one shows, by differentiating the period equations
(8), that the function
\[
S(u) = \frac{\partial^2 \sigma}{\partial u^2} + 4\eta_1\frac{\partial \sigma}{\partial \omega_1} + 4\eta_2\frac{\partial \sigma}{\partial \omega_2} + 4ru^2\sigma
\]
satisfies the period equations (8), and is therefore equal to $C\sigma(u)$,
where $C$ is independent of $u$. But the development of $\sigma(u)$ in % [sic] German source prints sigma(u); the argument requires S(u).
powers of $u$ begins with $u^3$, while $\sigma(u)$ begins with $u$,
and consequently we must have $C=0$.

By (7) the development of $\partial\sigma/\partial\omega_1$ and $\partial\sigma/\partial\omega_2$
begins only with $u^5$, and the first terms of $\partial^2\sigma/\partial u^2$ and $4ru^2\sigma$ are
$-g_2u^3/12$, $4ru^3$, from which $4r=g_2/12$ follows; accordingly
the differential equation for $\sigma$ takes the form:

\weq{10}{\frac{\partial^2 \sigma}{\partial u^2} + 4\eta_1\frac{\partial \sigma}{\partial \omega_1} + 4\eta_2\frac{\partial \sigma}{\partial \omega_2} + \frac{g_2}{12}u^2\sigma = 0,}

and incidentally there follow from (9) the relations:

\weq{11}{
\begin{aligned}
&\eta_1\frac{\partial \eta_1}{\partial \omega_1} + \eta_2\frac{\partial \eta_1}{\partial \omega_1} = -\frac{g_2}{48}\omega_1, \\ % [sic] The German source repeats partial eta_1 / partial omega_1 in the second term; the derivation requires partial eta_1 / partial omega_2.
&\eta_1\frac{\partial \eta_1}{\partial \omega_2} + \eta_2\frac{\partial \eta_2}{\partial \omega_2} = -\frac{g_2}{48}\omega_2. % [sic] The German source prints partial eta_1 / partial omega_2 in the first term; the derivation requires partial eta_2 / partial omega_1.
\end{aligned}
}

Now in the differential equation (10) the variables
$g_2, g_3$ are to be introduced in place of the variables $\omega_1, \omega_2$, and if we accordingly put
\[
\begin{aligned}
&\frac{\partial \sigma}{\partial \omega_1} = \frac{\partial \sigma}{\partial g_2}\frac{\partial g_2}{\partial \omega_1} + \frac{\partial \sigma}{\partial g_3}\frac{\partial g_3}{\partial \omega_1}, \\
&\frac{\partial \sigma}{\partial \omega_2} = \frac{\partial \sigma}{\partial g_2}\frac{\partial g_2}{\partial \omega_2} + \frac{\partial \sigma}{\partial g_3}\frac{\partial g_3}{\partial \omega_2}
\end{aligned}
\]
then (10) takes the form:

\weq{12}{\frac{\partial^2 \sigma}{\partial u^2} + h_2\frac{\partial \sigma}{\partial g_2} + h_3\frac{\partial \sigma}{\partial g_3} + \frac{g_2}{12}u^2\sigma = 0,}

where
\[
\begin{aligned}
h_2 &= 4\left(\eta_1\frac{\partial g_2}{\partial \omega_1} + \eta_2\frac{\partial g_2}{\partial \omega_2}\right), \\
h_3 &= 4\left(\eta_1\frac{\partial g_3}{\partial \omega_1} + \eta_2\frac{\partial g_3}{\partial \omega_2}\right).
\end{aligned}
\]

The quantities $h_2, h_3$, expressed through $g_2, g_3$, now follow
from the differential equation itself, if one substitutes the development
(7) for $\sigma$. Accordingly, up to the 7th power of $u$:
%CONT>>
